\NSPage{163}%
moment functions have the same \NSi{NS-F-P163-001}{O_m(N^{-1})} bounds by
integration, so \eqref{eq:C.16} remains valid inside the correction interval. Thus
the admissible stress cone persists throughout it.

At the correction interval's right endpoint
\NSi{NS-F-P163-002}{X_{\mathrm{rep}}}, the corrected profiles satisfy
\NSFormula{NS-F-P163-003}{display}{C.17}
\begin{equation}
  \widetilde{\mathfrak m}(X_{\mathrm{rep}},\eta)
    =\mathfrak m(X_{\mathrm{rep}},\eta),
  \qquad
  (\widetilde E,\widetilde U)=(E,U)
  \quad\text{for }X\geq X_{\mathrm{rep}}.
  \tag{C.17}\label{eq:C.17}
\end{equation}
Equality of the five radial moment functions then propagates by integration.
Lemma~\ref{lem:4.4} gives equality of the complete pressure and the functions
\NSi{NS-F-P163-004}{Q_s,N_s} for all
\NSi{NS-F-P163-005}{X\geq X_{\mathrm{rep}}}, including their parameter
derivatives. This equality preserves the terminal compensation already included
in the input moment vector \NSi{NS-F-P163-006}{\mathfrak m}. The exterior moments
and pressure normalized to vanish at radial infinity are consequently retained.
The first collar is unchanged because both modifications are supported later and
pressure was integrated from the unchanged axis datum. The last collar is
unchanged by the exact radial moment restoration. The two unused correction
patches retain their prescribed reference power law and
\NSi{NS-F-P163-007}{U=0}.

Finally fix one such finite \NSi{NS-F-P163-008}{N}. Every mixed radial and
parameter derivative of the resulting profiles is finite, with constants allowed
to depend on this \NSi{NS-F-P163-009}{N}. This choice is made before any limit in
the physical scale \NSi{NS-F-P163-010}{q} or any later dyadic band or correction
stage.
\end{proof}

\subsection{The two edges and the completed profile.}\label{sec:C.3}

In this subsection, we write \NSi{NS-F-P163-011}{E,U,\Pi} for the corrected
profiles of Proposition~\ref{prop:C.2} and form all associated quantities from
them. The interior cone inequalities and exact radial moments are now in place by
Proposition~\ref{prop:C.2}. In this subsection, we will prove
Proposition~\ref{prop:C.3} by verifying the remaining assertions at the annular
edges, where the stress \NSi{NS-F-P163-012}{\mathcal T_0} vanishes. Its direction
\NSi{NS-F-P163-013}{n} must extend smoothly and retain a strict cone margin, and
the derivatives of \NSi{NS-F-P163-014}{\mathcal T_0} must satisfy the common
exponential weight \NSi{NS-F-P163-015}{\zeta} in
Theorem~\ref{thm:4.6}. Both endpoint collars were preserved, so we can deduce
these properties from their original stress factorizations \eqref{eq:B.30} and
Equations~\eqref{eq:A.48} to~\eqref{eq:A.49}. In the weight below,
\NSi{NS-F-P163-016}{t_1>0} is the fixed activation width in
\NSi{NS-F-P163-017}{\log(X/X_a)} from Proposition~\ref{prop:4.10}.

\begin{proposition}\label{prop:C.3}
The profiles of Proposition~\ref{prop:C.2} satisfy all conclusions of
Theorem~\ref{thm:4.6}, with the weight
\NSFormula{NS-F-P163-018}{display}{C.18}
\begin{equation}
\begin{gathered}
  y_a=\log(X/X_a),
  \qquad y_b=\log(X_b/X),
  \qquad \delta=\min\{1,y_a,y_b\},\\
  \zeta=\exp\!\left(-\frac{t_1^2}{y_a^2}-\frac{4}{y_b^2}\right)
  \qquad (X_a<X<X_b),
\end{gathered}
\tag{C.18}\label{eq:C.18}
\end{equation}
extended by zero for \NSi{NS-F-P163-019}{X\leq X_a} and
\NSi{NS-F-P163-020}{X\geq X_b}.
\end{proposition}

\begin{proof}
The stress is zero up to \NSi{NS-F-P163-021}{X_a} by the analytic axis profile
and Proposition~\ref{prop:B.5}, and beyond \NSi{NS-F-P163-022}{X_b} by
Lemma~\ref{lem:A.8}. Proposition~\ref{prop:C.2} preserves these regions and makes
the admissible stress cone inequalities hold everywhere between them. If the
stress vanished at an interior point, the identity \eqref{eq:4.11} would give
\NSi{NS-F-P163-023}{p_s=(a,-b_s)}, hence
\NSi{NS-F-P163-024}{P_{\mathrm c}=v_s}, contradicting the strict quadratic cone test. This
proves the stress-support assertion in part~\textup{(ii)} of
Theorem~\ref{thm:4.6}; the unchanged axis profile also satisfies \eqref{eq:4.13}
for \NSi{NS-F-P163-025}{X\leq X_a}.

For the cone assertions in part~\textup{(iii)}, set
\NSi{NS-F-P163-026}{n=\mathcal T_0/|\mathcal T_0|} on
\NSi{NS-F-P163-027}{X_a<X<X_b}. The bounds
\NSi{NS-F-P163-028}{F>0}, \NSi{NS-F-P163-029}{a>0}, and
\NSi{NS-F-P163-030}{v_s>2} follow in the interior from the positive profile
construction and the admissible stress cone. At the inner endpoint
\NSi{NS-F-P163-031}{a=p_{1,\mathrm r}>0}, while Proposition~\ref{prop:B.5} gives
\NSi{NS-F-P163-032}{v_s>2+c} there and the analytic axis profile has
\NSi{NS-F-P163-033}{F>0}. At the outer endpoint they follow from
\eqref{eq:A.56}, with \NSi{NS-F-P163-034}{b_s=0}, and positivity of the heat
factor. Continuity and compactness give the asserted positive lower bounds on the
closed annulus.

The inner edge has the preserved factorization
\NSi{NS-F-P163-035}{\mathcal T_0=e_aB_0} in \eqref{eq:B.30}, with
\NSi{NS-F-P163-036}{e_a=e^{-t_1^2/y_a^2}g_a(y_a)},
\NSi{NS-F-P163-037}{g_a(0)>0}, and
\NSi{NS-F-P163-038}{B_0(0,\eta)=Fp_{s,\mathrm r}\neq0}. Its direction is therefore
smooth, and at the edge it is a

\NSPage{164}%
positive scalar multiple of the limiting shear
\NSi{NS-F-P164-001}{(a,-b_s)=a(1,t_s)}. In particular
\NSi{NS-F-P164-002}{n_z-t_sn_\theta=0} and
\NSi{NS-F-P164-003}{n_\theta+t_sn_z>0} there. At the outer edge, the preserved
factors of Proposition~\ref{prop:A.10} give
\NSFormula{NS-F-P164-004}{display}{none}
\[
  n=\frac{(b_\theta,y_b^6b_z)}
          {\sqrt{b_\theta^2+y_b^{12}b_z^2}},
  \qquad b_\theta(0,\eta)>0.
\]
Thus the outer limiting direction is \NSi{NS-F-P164-005}{(1,0)}, and the
limiting shear ratio is \NSi{NS-F-P164-006}{t_s=0}. All these statements are
smooth through \NSi{NS-F-P164-007}{\eta=\pm1}; the heat factor is evaluated only
on its nonnegative argument domain.

The edge limits allow the interior cone inequalities to be made uniform on the
closed annulus. In the interior, \eqref{eq:4.11} implies
\NSFormula{NS-F-P164-008}{display}{none}
\[
  n_\theta+t_sn_z
    =\frac{P_{\mathrm c}-v_s}{|p_s-(a,-b_s)|}>0,
  \qquad
  n_z-t_sn_\theta
    =\frac{J_{\mathrm c}}{|p_s-(a,-b_s)|}.
\]
The admissible stress cone makes
\NSi{NS-F-P164-009}{2-(v_s-2)(n_z-t_sn_\theta)^2/(n_\theta+t_sn_z)^2}
positive there. Both this expression and its positive denominator extend
continuously to the two edges; the former has value two at either edge by the
factorizations just proved. Their compact positive minima yield, for a single
\NSi{NS-F-P164-010}{0<\kappa<2},
\NSFormula{NS-F-P164-011}{display}{none}
\[
  n_\theta+t_sn_z\geq\kappa,
  \qquad
  (v_s-2)(n_z-t_sn_\theta)^2
    \leq(2-\kappa)(n_\theta+t_sn_z)^2.
\]
Together with the endpoint directions and positive lower bounds above, this
proves part~\textup{(iii)} of Theorem~\ref{thm:4.6}.

For part~\textup{(iv)}, it remains to prove the weighted estimates
\NSFormula{NS-F-P164-012}{display}{C.19}
\begin{equation}
  |\mathcal T_0|\geq c\zeta,
  \qquad
  |\partial^I\mathcal T_0|\leq C_I\zeta\delta^{-m_I}
  \qquad (X_a<X<X_b,\ -1\leq\eta\leq1)
  \tag{C.19}\label{eq:C.19}
\end{equation}
for every fixed mixed profile derivative
\NSi{NS-F-P164-013}{\partial^I}, with finite constants independent of physical
\NSi{NS-F-P164-014}{q}.

On an inner collar \NSi{NS-F-P164-015}{\zeta} differs from
\NSi{NS-F-P164-016}{e^{-t_1^2/y_a^2}} by a smooth positive factor bounded above
and below; the bounds \eqref{eq:B.31} therefore imply \eqref{eq:C.19} there. On
an outer collar it differs from \NSi{NS-F-P164-017}{e^{-4/y_b^2}} in the same
way. The angular factor
\NSi{NS-F-P164-018}{e^{-4/y_b^2}y_b^{-3}b_\theta}, with
\NSi{NS-F-P164-019}{b_\theta} bounded below, proves the lower bound, and
\eqref{eq:A.51} proves every fixed derivative upper bound. On the remaining
compact interior interval, both \NSi{NS-F-P164-020}{\zeta} and
\NSi{NS-F-P164-021}{|\mathcal T_0|} have positive minima and all
fixed derivatives are bounded. Combining the three regions proves
\eqref{eq:C.19}. The exponential factors also show that extension of the stress
by zero across both edges is smooth. The definition \eqref{eq:C.18} also makes
\NSi{NS-F-P164-022}{\zeta} positive in the interior and flat at both edges,
completing part~\textup{(iv)} of Theorem~\ref{thm:4.6}.

For the regularity in part~\textup{(i)}, Proposition~\ref{prop:B.2} supplies the
smooth axis profiles with \NSi{NS-F-P164-023}{F>0}, and
Corollary~\ref{cor:B.6} preserves the common analytic rectangle
\NSi{NS-F-P164-024}{[0,X_{\mathrm{an}}]\times[-1,1]}. The factors give Cartesian
smoothness across the axis by Equations~\eqref{eq:4.4} to~\eqref{eq:4.5}. The
later modifications in Proposition~\ref{prop:C.2} are smooth and preserve
positivity, with finite mixed derivative bounds on each fixed radial interval.
The inner directional margin is part~\textup{(iii)}, proved above. The heat
profile is smooth through \NSi{NS-F-P164-025}{\eta=\pm1} by
Lemma~\ref{lem:A.6}; no analyticity of its factor at zero is needed. The pressure
normalization \eqref{eq:4.25} follows from Lemma~\ref{lem:A.5} and
Propositions~\ref{prop:A.7} and~\ref{prop:B.8} and is retained by
\eqref{eq:C.17}. This completes part~\textup{(i)}. Differentiating this pressure
formula gives \NSi{NS-F-P164-026}{\Pi_X=E^2/(2X)=F^2}, with its smooth extension
at \NSi{NS-F-P164-027}{X=0}. The exact tangential residual identities follow
from Proposition~\ref{prop:4.2}; together with the support and axis checks above,
they complete part~\textup{(ii)}.

For part~\textup{(v)}, the exact exterior moments \eqref{eq:4.28} are supplied by
Propositions~\ref{prop:A.7} and~\ref{prop:B.8} and retained by
\eqref{eq:C.17}. The same restoration preserves the heat exterior
\eqref{eq:4.29}, whose exact heat evolution and smoothness are proved in
Lemma~\ref{lem:A.6}. We take
\NSi{NS-F-P164-028}{X_{\mathrm v}=X_{\mathrm{end}}\in(X_a,X_b)},
